\documentclass[12pt]{beamer}
\usepackage{ceai-beamer}
\usepackage{ragged2e}
\AtBeginEnvironment{frame}{\justifying}

\title{CE 488/5588: Applied AI}
\subtitle{Topic 10: Generalization and Model Selection}
\author{}
\date{}

\begin{document}

\begin{frame}
  \titlepage
\end{frame}

\begin{frame}{Outline}
  \tableofcontents
\end{frame}

% ==========================================================
% Generalization
% ==========================================================
\section{Generalization}

\begin{frame}{Why this topic matters}
  \begin{conceptbox}{Core question}
    The real objective of machine learning is not to score well on the data we already saw,
    but to perform reliably on unseen data drawn from the same or similar conditions.
  \end{conceptbox}
  \begin{itemize}
    \item Apparent training or test success can be misleading when the split is small or lucky.
    \item In engineering practice, labeled data are often expensive, limited, and noisy.
    \item This makes \textbf{generalization} the central criterion for judging model quality.
  \end{itemize}
\end{frame}

\begin{frame}{Learning objectives}
  \begin{itemize}
    \item Explain what generalization means in supervised learning.
    \item Distinguish underfitting and overfitting.
    \item Understand why bias and variance pull in opposite directions.
    \item Use re-sampling methods to estimate model performance more robustly.
    \item Describe model selection through threshold tuning, hyperparameter tuning, and model comparison.
    \item Summarize why ensemble methods often improve robustness.
  \end{itemize}
\end{frame}

\begin{frame}{A simple train-test split is not the whole story}
  \begin{itemize}
    \item Typical supervised workflow:
      \begin{itemize}
        \item train on $\mathcal{D}_{tr}$,
        \item evaluate on $\mathcal{D}_{te}$.
      \end{itemize}
    \item This is necessary, but one split alone may not be trustworthy.
    \item With finite data,
      \[
        \mathcal{D}=\{(\bm{x}_i,c_i)\mid i=1,\ldots,N\}
      \]
      must be partitioned, and different partitions can produce noticeably different estimates.
    \item A model that changes quality drastically under a different split has weak generalization.
  \end{itemize}
\end{frame}

\begin{frame}{What is generalization?}
  \begin{conceptbox}{Working definition}
    Generalization is a model's ability to capture underlying patterns so that it remains effective
    on unseen samples rather than merely memorizing the training set.
  \end{conceptbox}
  \begin{itemize}
    \item Good generalization means low error on future data from the same task distribution.
    \item Poor generalization usually appears in two opposite failure modes:
      \begin{itemize}
        \item \textbf{Underfitting}: the model is too simple to capture the structure.
        \item \textbf{Overfitting}: the model learns noise, accidental detail, or sampling quirks.
      \end{itemize}
  \end{itemize}
  \bhref{https://see.stanford.edu/Course/CS229/53}{Stanford bias-variance}
\end{frame}

\begin{frame}{Underfitting vs. overfitting}
  \begin{columns}[T]
    \column{0.48\textwidth}
      \begin{warningbox}{Underfitting}
        \begin{itemize}
          \item Model capacity is too low.
          \item Training error stays high.
          \item Test error is also high.
          \item Important signal is missed.
        \end{itemize}
      \end{warningbox}
    \column{0.48\textwidth}
      \begin{warningbox}{Overfitting}
        \begin{itemize}
          \item Model capacity is too high for the available data.
          \item Training error becomes very low.
          \item Test performance deteriorates.
          \item Noise is treated as signal.
        \end{itemize}
      \end{warningbox}
  \end{columns}
\end{frame}

\begin{frame}{Generalization as a tradeoff}
  \begin{itemize}
    \item Bias and variance offer a useful lens for understanding model behavior.
    \item \textbf{Bias} is error caused by overly restrictive modeling assumptions.
      \begin{itemize}
        \item High bias tends to produce underfitting.
      \end{itemize}
    \item \textbf{Variance} is error caused by excessive sensitivity to the sampled training data.
      \begin{itemize}
        \item High variance tends to produce overfitting.
      \end{itemize}
    \item Good generalization requires a balance rather than maximizing complexity.
  \end{itemize}
\end{frame}

\begin{frame}{Bias-variance intuition}
\vspace{-30pt}
  \begin{figure}
    \centering
    \includegraphics[width=0.36\textwidth]{bias-variance-01.pdf}\hfill
    \includegraphics[width=0.45\textwidth]{bias-variance-02.pdf}
  \end{figure}
  \begin{itemize}
    \item As complexity increases, bias usually decreases.
    \item Past a certain point, variance rises quickly and hurts unseen-data performance.
    \item The best model is usually near the minimum of validation or test error, not training error.
  \end{itemize}
\end{frame}

\begin{frame}{How we improve generalization in practice}
  \begin{itemize}
    \item Four practical levers appear repeatedly:
      \begin{itemize}
        \item \textbf{Regularization}: constrain model complexity.
        \item \textbf{Re-sampling}: estimate performance more robustly from limited data.
        \item \textbf{Model selection}: choose thresholds, hyperparameters, or model families wisely.
        \item \textbf{Ensembling}: combine multiple learners to reduce instability.
      \end{itemize}
    \item This topic emphasizes re-sampling, selection, and ensemble strategies.
  \end{itemize}
\end{frame}

% ==========================================================
% Re-sampling Methods
% ==========================================================
\section{Re-sampling Methods}

\begin{frame}{Why re-sampling is needed}
  \begin{itemize}
    \item A single split wastes information when data are scarce.
    \item Re-sampling reuses the available dataset to estimate performance more reliably.
    \item It supports three practical goals:
      \begin{itemize}
        \item more robust performance measurement,
        \item better detection of overfitting or underfitting,
        \item more defensible model selection.
      \end{itemize}
  \end{itemize}
  \bhref{https://scikit-learn.org/stable/modules/cross_validation.html}{sklearn cross-validation}
\end{frame}

\begin{frame}{K-fold cross-validation: idea}
  \begin{conceptbox}{Main idea}
    Partition the dataset into $K$ folds, train on $K-1$ folds, test on the remaining fold,
    and repeat until every fold has served once as the validation fold.
  \end{conceptbox}
  \begin{itemize}
    \item Every sample is used for both training and validation, but not in the same iteration.
    \item Typical choices are $K=5$ or $K=10$.
    \item The reported score is usually the mean across folds, often with standard deviation.
  \end{itemize}
\end{frame}

\begin{frame}{K-fold cross-validation workflow}
  \begin{enumerate}
    \item Split $\mathcal{D}$ into $K$ approximately equal folds.
    \item For fold $j=1,\ldots,K$:
      \begin{itemize}
        \item train on the other $K-1$ folds,
        \item evaluate on fold $j$.
      \end{itemize}
    \item Aggregate the scores:
      \[
        \bar{s}=\frac{1}{K}\sum_{j=1}^{K}s_j
      \]
    \item Prefer models with strong mean performance and acceptable variability across folds.
  \end{enumerate}
\end{frame}

\begin{frame}{Why K-fold cross-validation helps}
  \begin{itemize}
    \item It reduces dependence on one arbitrary holdout split.
    \item It exposes instability:
      \begin{itemize}
        \item large variance across folds suggests sensitivity to data selection,
        \item consistently weak scores suggest underfitting or poor features.
      \end{itemize}
    \item It is widely used for hyperparameter tuning and model comparison.
    \item Stratified variants are often preferred for imbalanced classification.
  \end{itemize}
\end{frame}

\begin{frame}{Bootstrapping: idea}
  \begin{conceptbox}{Sampling with replacement}
    Bootstrapping repeatedly draws new datasets from the original dataset by sampling with replacement.
  \end{conceptbox}
  \begin{itemize}
    \item Each bootstrap sample has the same nominal size as the original dataset.
    \item Some original samples appear multiple times; some do not appear at all.
    \item Repeated bootstrap fits can estimate variability in model behavior or statistics.
  \end{itemize}
  \bhref{https://www.stat.cmu.edu/~larry/=sml/Boot.pdf}{CMU bootstrap notes}
\end{frame}

\begin{frame}{Cross-validation vs. bootstrapping}
  \scriptsize
  \renewcommand{\arraystretch}{1.25}
  \setlength{\tabcolsep}{4pt}
  \begin{tabular}{|p{2.1cm}|p{3.2cm}|p{3.2cm}|}
    \hline
    Method & Strengths & Caveats \\
    \hline
    K-fold CV & Non-overlapping validation folds; natural for model comparison; widely used in ML workflows & Can still be noisy on very small datasets; repeated training cost grows with $K$ \\
    \hline
    Bootstrap & Useful for estimating variability and uncertainty from limited data; reuses samples aggressively & Samples repeat with replacement; interpretation differs from fold-based validation \\
    \hline
  \end{tabular}
  \vspace{0.4em}
  \begin{itemize}
    \item Cross-validation is usually the default for model evaluation and tuning.
    \item Bootstrapping is especially useful when distributional uncertainty is of interest.
  \end{itemize}
\end{frame}

\begin{frame}{Practical cautions for re-sampling}
  \begin{warningbox}{Avoid leakage}
    Any preprocessing learned from data -- scaling, feature selection, imputation, threshold tuning --
    must be fit inside each training split, not on the full dataset ahead of time.
  \end{warningbox}
  \begin{itemize}
    \item Leakage produces over-optimistic performance estimates.
    \item Time-series, grouped, or spatial data often require specialized split strategies.
    \item Re-sampling estimates model quality; it does not magically fix poor data quality.
  \end{itemize}
\end{frame}

% ==========================================================
% Model Selection
% ==========================================================
\section{Model Selection}

\begin{frame}{Three model-selection scenarios}
  \begin{itemize}
    \item \textbf{Threshold selection}: choose the decision cutoff for turning scores into labels.
    \item \textbf{Hyperparameter tuning}: choose values such as $C$, kernel width, depth, or layer count.
    \item \textbf{Architecture selection}: choose among different model families or complexity levels.
  \end{itemize}
  \begin{conceptbox}{Common principle}
    Every selection decision should be based on expected unseen-data performance rather than training fit alone.
  \end{conceptbox}
\end{frame}

\begin{frame}{Classification threshold selection}
  \begin{itemize}
    \item Many classifiers produce a score or estimated probability $P(c\mid\bm{x})$.
    \item A hard decision is then made by thresholding that score.
    \item Default binary rule:
      \[
        c =
        \begin{cases}
          1 & P(c\mid\bm{x}) \geq 0.5 \\
          0 & P(c\mid\bm{x}) < 0.5
        \end{cases}
      \]
    \item A threshold of $0.5$ is a convention, not a law.
  \end{itemize}
  \bhref{https://sklearn.org/stable/modules/classification_threshold.html}{sklearn threshold tuning}
\end{frame}

\begin{frame}{Why the threshold should depend on the application}
  \begin{itemize}
    \item Different applications attach different costs to false positives and false negatives.
    \item Example: medical diagnosis
      \begin{itemize}
        \item false negative: miss a dangerous condition,
        \item false positive: trigger follow-up testing or anxiety.
      \end{itemize}
    \item If missing a positive case is more severe, lower the threshold to increase recall.
    \item If false alarms are very expensive, raise the threshold to improve precision.
  \end{itemize}
\end{frame}

\begin{frame}{ROC curves for threshold analysis}
  \begin{itemize}
    \item ROC tracks performance over all thresholds using:
      \[
        TPR=\frac{TP}{TP+FN}, \qquad FPR=\frac{FP}{FP+TN}.
      \]
    \item Each threshold gives one operating point.
    \item Curves closer to the upper-left corner indicate better ranking behavior.
    \item Area under the ROC curve (AUC) summarizes threshold-wide separability.
  \end{itemize}
  \bhref{https://developers.google.com/machine-learning/crash-course/classification/roc-and-auc}{Google ROC/AUC}
\end{frame}

\begin{frame}{ROC visualization and reading}
  \begin{figure}
    \centering
    \includegraphics[width=0.72\textwidth]{Topic08fig01.png}
    \caption{ROC curve with threshold-dependent tradeoff and baseline.}
  \end{figure}
  \begin{itemize}
    \item $AUC=1$ means perfect ranking; $AUC=0.5$ is random ranking.
    \item ROC is useful when class distributions are reasonably balanced.
  \end{itemize}
\end{frame}

\begin{frame}{Precision-recall curves}
  \begin{itemize}
    \item PR curves focus directly on the positive class:
      \[
        Precision=\frac{TP}{TP+FP}, \qquad Recall=\frac{TP}{TP+FN}.
      \]
    \item They are especially informative when positive cases are rare.
    \item A strong model keeps precision high while recall increases.
    \item Use PR together with ROC when class imbalance matters.
  \end{itemize}
  \bhref{https://scikit-learn.org/stable/modules/generated/sklearn.metrics.precision_recall_curve.html}{sklearn PR curve}
\end{frame}

\begin{frame}{PR visualization and imbalance warning}
  \begin{columns}[T]
    \column{0.58\textwidth}
      \begin{figure}
        \centering
        \includegraphics[width=\textwidth]{Topic08fig02.png}
      \end{figure}
    \column{0.45\textwidth}
      \begin{warningbox}{Why PR matters}
        \begin{itemize}
          \item ROC can look strong when true negatives dominate.
          \item PR exposes whether the retrieved positives are actually useful.
          \item For rare-event detection, PR is often the more revealing curve.
        \end{itemize}
      \end{warningbox}
  \end{columns}
\end{frame}

\begin{frame}{Hyperparameter tuning}
  \begin{itemize}
    \item Hyperparameters are design choices set by the user rather than learned directly from data.
    \item Examples:
      \begin{itemize}
        \item logistic regression or SVM regularization parameter $C$,
        \item kernel type and RBF width $\gamma$,
        \item tree depth, number of estimators, network depth.
      \end{itemize}
    \item Hyperparameter tuning is a model-selection problem because different settings change model capacity.
  \end{itemize}
\end{frame}

\begin{frame}{A practical tuning workflow}
  \begin{enumerate}
    \item Choose a metric aligned with the application goal.
    \item Define a candidate grid or search space for hyperparameters.
    \item Evaluate each setting with cross-validation.
    \item Compare mean score and variability across folds.
    \item Retrain the selected setting on the full training set.
    \item Report final performance on a held-out test set.
  \end{enumerate}
  \bhref{https://sklearn.org/stable/model_selection.html}{sklearn model selection}
\end{frame}

\begin{frame}{Choosing among different model architectures}
  \begin{itemize}
    \item Sometimes the decision is not just a parameter value but a whole model family:
      \begin{itemize}
        \item MLP vs. SVM,
        \item decision tree vs. logistic regression,
        \item shallow vs. deep models.
      \end{itemize}
    \item More flexible models can capture richer patterns but usually demand more data and stronger regularization.
    \item Simpler models are easier to interpret and often more stable on small datasets.
  \end{itemize}
\end{frame}

\begin{frame}{Occam's razor and information criteria}
  \begin{conceptbox}{Parsimony principle}
    If multiple models explain the data similarly well, prefer the simpler one unless extra complexity
    clearly improves expected performance.
  \end{conceptbox}
  \begin{itemize}
    \item Probability-based criteria penalize complexity while rewarding fit.
    \item Common examples: AIC, BIC, MDL, and Structural Risk Minimization (SRM).
    \item In this course, these are awareness concepts rather than implementation details.
  \end{itemize}
  \bhref{https://sta210-s22.github.io/website/supplemental/model-selection-criteria.html}{AIC/BIC notes}
\end{frame}

% ==========================================================
% Ensemble Models
% ==========================================================
\section{Ensemble Models}

\begin{frame}{Why combine models?}
  \begin{itemize}
    \item Different models make different mistakes.
    \item Combining them can reduce variance, improve robustness, and sometimes improve accuracy.
    \item Ensemble thinking is especially attractive when no single architecture dominates consistently.
  \end{itemize}
  \begin{conceptbox}{Key intuition}
    Weak or unstable learners can become stronger when aggregated appropriately.
  \end{conceptbox}
\end{frame}

\begin{frame}{Three major ensemble strategies}
  \begin{itemize}
    \item \textbf{Bagging}
      \begin{itemize}
        \item Train multiple models on different sampled subsets.
        \item Aggregate by voting or averaging.
        \item Main effect: variance reduction.
      \end{itemize}
    \item \textbf{Boosting}
      \begin{itemize}
        \item Train models sequentially, emphasizing previous mistakes.
        \item Main effect: strong learner built from focused weak learners.
      \end{itemize}
    \item \textbf{Stacking}
      \begin{itemize}
        \item Train diverse base models, then learn a meta-model to combine them.
      \end{itemize}
  \end{itemize}
\end{frame}

\begin{frame}{Bagging, boosting, and stacking compared}
  \scriptsize
  \renewcommand{\arraystretch}{1.2}
  \setlength{\tabcolsep}{4pt}
  \begin{tabular}{|p{1.8cm}|p{2.5cm}|p{2.7cm}|p{2.2cm}|}
    \hline
    Method & Training pattern & Main benefit & Typical example \\
    \hline
    Bagging & Parallel models on re-sampled data & Lower variance, better stability & Random forest \\
    \hline
    Boosting & Sequential learners focusing on errors & Stronger predictive accuracy & AdaBoost, gradient boosting \\
    \hline
    Stacking & Diverse base learners plus meta-learner & Learns how to combine strengths & Blended heterogeneous models \\
    \hline
  \end{tabular}
  \vspace{0.4em}
  \bhref{https://scikit-learn.org/stable/modules/ensemble.html}{sklearn ensembles}
\end{frame}

\begin{frame}{When ensembles help most}
  \begin{itemize}
    \item Complex datasets where no single model captures every pattern well.
    \item High-variance learners whose predictions fluctuate with the sampled data.
    \item Situations where robustness matters more than having one elegant model.
  \end{itemize}
  \begin{warningbox}{But not free}
    \begin{itemize}
      \item Ensembles increase computation and deployment complexity.
      \item Interpretability often becomes harder.
      \item If all component models make the same mistake, ensembling adds little value.
    \end{itemize}
  \end{warningbox}
\end{frame}

\begin{frame}{Topic 10 big picture}
  \begin{itemize}
    \item Generalization is the ultimate target.
    \item Re-sampling provides better evidence about expected unseen-data performance.
    \item Model selection chooses the operating point, the hyperparameters, and sometimes the architecture.
    \item Ensemble learning is a practical hedge against instability and single-model weakness.
  \end{itemize}
\end{frame}

% ==========================================================
% Conclusion Remarks
% ==========================================================
\section{Conclusion Remarks}

\begin{frame}{Key takeaways}
  \begin{itemize}
    \item Strong training fit does not guarantee strong deployment behavior.
    \item Underfitting and overfitting are best understood through bias-variance balance.
    \item K-fold cross-validation and bootstrapping improve evaluation with limited data.
    \item Threshold tuning, hyperparameter tuning, and architecture selection are all model-selection tasks.
    \item Ensemble methods often improve robustness by combining complementary models.
  \end{itemize}
\end{frame}

\begin{frame}{Suggested workflow for practice}
  \begin{enumerate}
    \item Start with a simple baseline model.
    \item Pick a metric tied to the application cost function.
    \item Use cross-validation for tuning and comparison.
    \item Inspect ROC and PR behavior when threshold choice matters.
    \item Evaluate whether a simpler model is good enough before moving to an ensemble.
  \end{enumerate}
\end{frame}

\begin{frame}{References for further study}
  \begin{itemize}
    \item \bhref{https://see.stanford.edu/Course/CS229/53}{Stanford: Bias-Variance Tradeoff}
    \item \bhref{https://scikit-learn.org/stable/modules/cross_validation.html}{scikit-learn: Cross-Validation}
    \item \bhref{https://www.stat.cmu.edu/~larry/=sml/Boot.pdf}{CMU: Bootstrap Chapter}
    \item \bhref{https://sklearn.org/stable/modules/classification_threshold.html}{scikit-learn: Threshold Tuning}
    \item \bhref{https://developers.google.com/machine-learning/crash-course/classification/roc-and-auc}{Google: ROC and AUC}
    \item \bhref{https://scikit-learn.org/stable/modules/generated/sklearn.metrics.precision_recall_curve.html}{scikit-learn: PR Curves}
    \item \bhref{https://sklearn.org/stable/model_selection.html}{scikit-learn: Model Selection}
    \item \bhref{https://scikit-learn.org/stable/modules/ensemble.html}{scikit-learn: Ensemble Methods}
  \end{itemize}
\end{frame}

\begin{frame}{Q\&A}
  \centering
  Questions?
\end{frame}

\end{document}
